\documentclass[11pt]{article}
\usepackage[margin=1in]{geometry}
\usepackage{amsmath,amssymb,booktabs,graphicx,microtype,tabularx,xcolor}
\usepackage[colorlinks=true,linkcolor=blue!50!black,citecolor=blue!50!black]{hyperref}
\usepackage{caption}
\captionsetup{font=small}
\usepackage[round]{natbib}
\newcommand{\E}{\mathbb{E}}
\newcommand{\norm}[1]{\lVert #1 \rVert}
\title{\textbf{Running the model forward, and measuring how far it lands}\\[2pt]
\large The sampling procedure and the energy-distance yardstick behind the\\
generative-failure figure, with the numbers from the fixed engine (R36i)}
\author{ProductGPT technical note (R36)}
\date{19 October 2026; revised 31 October 2026}

\begin{document}
\maketitle

\begin{abstract}
\noindent The generative-failure figure makes two claims. A model that predicts
the next decision as well as any other family produces whole purchase histories
far from the real ones when it is run on its own output. And most of that
distance comes from feeding the model its own decisions and acquisitions, not
from the per-step distribution or from the randomness of the loot boxes. This
note documents how each generator is sampled
(Sections~\ref{sec:objects}--\ref{sec:generators}), how the distance between
generated and real histories is defined and calibrated
(Section~\ref{sec:metric}), and the numbers from R36i, which re-ran every arm
on the fixed rollout engine with one set of 300 customers and one band
(Section~\ref{sec:results}). On the fixed engine, free running lands at 41
times the band and loses to independent draws. Scheduled sampling closes 59\%
of the gap, and the Markov chain stays closer than every neural generator.
Teacher forcing is outside the band (5.3 times), so only the weaker claim
survives: most of the distance comes from the feedback loop. Arm C and teacher
forcing read the realised holdout; they diagnose the failure but cannot
simulate a policy. Section~\ref{sec:weakness} shows why energy distance is
nearly blind to a wrong dependence structure, and Section~\ref{sec:perstat}
breaks the distance down by statistic. Later runs on the same yardstick show
classical parametric choice models generating ten times closer than free-running
ProductGPT while predicting worse. Section~\ref{sec:caveats} lists the
qualifications.
\end{abstract}

\section{What has to be built}\label{sec:question}

Everything up to R35 evaluated the model one step at a time: given a customer's
\emph{true} history, how well is the next decision predicted? A marketing
application wants more than that. To simulate a policy, the model has to
generate whole sequences, conditioning on what it generated a moment earlier.
If the model's conditionals were correctly specified, the chain rule would
deliver valid sequences for free. The two criteria come apart only under
misspecification, and they come apart for a specific reason. Training evaluates
the model on histories drawn from the \emph{data}; a rollout visits histories
drawn from the \emph{model}. A sequence-level failure alongside good one-step
fit is therefore a specification test that the likelihood cannot perform.

Testing this needs two instruments. The first is a way to generate holdout
histories under different feedback regimes, so the source of any failure can
be located. The second is a distance between a \emph{set} of generated histories
and the set of real ones, with a reference scale that says how large a distance
real data produces against itself.

\section{Objects and notation}\label{sec:objects}

Customers are indexed $i = 1,\dots,N$, and their occasions (rows of the panel)
$t = 1,\dots,T_i$, in calendar order. Each occasion belongs to a campaign
$c(t)$. The holdout is $\mathcal{H}_i = \{t : c(t) \ge 28\}$, campaigns 28--30,
and $t_i^0 = \min \mathcal{H}_i$ is the customer's first holdout occasion.

\begin{itemize}
\item \textbf{Decision} $y_{it} \in \{1,\dots,9\}$: one or ten pulls
  (Buy1, Buy10) on one of four banners (regular, figure A, figure B, weapon),
  or NotBuy.
\item \textbf{Offers} $o_{it}$: the featured products on each banner at $t$.
  Offers are set by the firm and are treated as exogenous.
\item \textbf{Acquisitions} $a_{it}$: up to ten product identifiers received
  as a result of $y_{it}$.
\item \textbf{Elapsed time} $\Delta_{it}$ since the previous occasion.
\end{itemize}

\noindent Row $t$ of the model's input carries the outcome of the decision made
at $t-1$. The model's information at $t$ is therefore
\begin{equation}
x_{it} = \big(o_{i,1:t},\; a_{i,0:t-1},\; y_{i,0:t-1},\; \Delta_{i,1:t}\big),
\label{eq:info}
\end{equation}
truncated to the last 384 rows, plus a summary of the inventory accumulated
before that window. Everything the model knows about holdings, duplicates and
satiation is a function of the acquisition stream $a$. The trained model is a
nine-way conditional distribution $p_\theta(y \mid x)$, fitted by one-step
maximum likelihood on observed histories.

What a decision yields is determined outside the model, by an acquisition
environment
\begin{equation}
a_{it} \sim Q\big(\,\cdot \mid y_{it}, o_{it}, \pi_{it}\big),
\qquad \pi_{i,t+1} = U(\pi_{it}, y_{it}, a_{it}),
\label{eq:env}
\end{equation}
where $\pi_{it}$ is the customer's pity state (pulls since the last 4-star and
5-star, per banner) and $U$ is its deterministic update.

\section{The acquisition environment \texorpdfstring{$Q$}{Q}}\label{sec:env}

\paragraph{How many items.} A purchase fills $n$ slots, with $n$ drawn from the
empirical distribution of filled slots given the decision. This is calibrated
rather than fixed at 1 or 10 because the event-level attribution in the data is
imperfect: a Buy10 fills ten slots only 72\% of the time, and one slot 23\% of
the time.

\paragraph{Rarity of each pull.} Rarity follows the published rates with pity.
If $s$ pulls have passed since the last 5-star on a banner, the next pull is a
5-star with probability
\begin{equation}
p_5(s) =
\begin{cases}
\beta & s+1 < \text{soft},\\
\min\{1,\; \beta + (s+2-\text{soft})\cdot 10\beta\} & \text{soft} \le s+1 < \text{hard},\\
1 & s+1 \ge \text{hard}.
\end{cases}
\end{equation}
4-stars follow the same rule with their own counter. The parameters are those
published for the game:

\begin{center}
\small
\begin{tabular}{lcccccc}
\toprule
banner & $\beta_5$ & soft$_5$ & hard$_5$ & $\beta_4$ & soft$_4$ & hard$_4$\\
\midrule
regular, figure A, figure B & 0.006 & 74 & 90 & 0.051 & 9 & 10\\
weapon & 0.007 & 63 & 80 & 0.060 & 9 & 10\\
\bottomrule
\end{tabular}
\end{center}

\noindent The data agree with them closely: measured 5-star rates on the
calibration campaigns are 1.25\% (figure banners) and 1.73\% (weapon), against
published effective rates of about 1.6\% and 1.85\%. The 4-star rate is 13\%,
against a published 13\%.

\paragraph{Which product.} When a 5-star lands on a banner whose featured slot
is filled, it is the featured product, a standard-pool product, or another
limited product. The three shares are taken from the calibration campaigns. A
non-featured 5-star, and every 3- or 4-star, is drawn by its calibration
frequency among the products that banner actually yielded at that rarity.

The game's ``50/50 guarantee'' (after losing the coin flip once, the next 5-star
is featured; two losses on the weapon banner) is implemented, but it is switched
off by the data. Under a guarantee after $L$ losses, the long-run featured share
is $1/\sum_{s=0}^{L}(1-q)^s$. That share can never fall below 0.5 on a character
banner or 0.333 on the weapon banner. The measured shares are 0.323 (figure A),
0.146 (figure B) and 0.335 (weapon), so no $q$ reproduces them, and the
environment uses the unconditional three-way split (R36h). This is a finding
about the data's attribution of acquisitions to banners, not about the game.

\paragraph{What it is fitted on.} Calibration rows ($c(t) < 28$) of the first
600 customers in the \emph{training} set. The environment never sees the
customers being scored or any holdout campaign.

\paragraph{Warm-up.} Before a customer's rollout starts, the pity counters are
advanced by replaying that customer's real history up to $t_i^0$, so each
simulated customer enters the holdout with the pity state the real customer had.

\section{The generators}\label{sec:generators}

\subsection{Common protocol}

\begin{enumerate}
\item \textbf{Customers.} $N=300$ out-of-sample customers, drawn at random
  (seed 0) from the cell that produces the headline NLL (customers held out of
  training, scored on the holdout campaigns): 22,490 holdout occasions in all.
\item \textbf{Models and runs.} The neural arms use b33, the gen-5 model with
  campaign effects and a holding half-life of 2, and b35, the same model
  trained with scheduled sampling. Every arm has three runs: training seeds
  1--3 for model arms, three draw seeds otherwise.
\item \textbf{Context window.} The attention context is cut to the last 384
  occasions to fit in memory, with everything acquired before the window carried
  in. This costs b33 $+0.0007$ nats of teacher-forced NLL against the full 1,024
  (0.8615 against 0.8608 on the same 300 customers).
\item \textbf{Fixed occasion grid.} The occasions and their elapsed times are
  held at their observed values; only the decision at each occasion is
  generated. The model consumes elapsed time but does not predict it, so it
  cannot generate the grid. What is tested is therefore the composition and
  dependence of decisions \emph{given when occasions arrive}, not purchase
  timing or frequency.
\item \textbf{Real past, real offers.} Rows before $t_i^0$ are the customer's
  real history. Offers are real throughout.
\item \textbf{Replicates.} $M=8$ independent replicates per customer, so 2,400
  generated histories are compared with 300 real ones.
\item \textbf{Sampling rule.} Ancestral sampling at temperature 1 with no top-$k$
  or nucleus truncation: draw $u \sim U(0,1)$ and take the first class whose
  cumulative probability exceeds $u$. Truncated decoding would change the very
  distribution under test.
\end{enumerate}

\subsection{Free running}

For each holdout occasion $t \in \mathcal{H}_i$, in order:
\begin{align}
\tilde y_{it} &\sim p_\theta\big(\,\cdot \mid \tilde x_{it}\big), \\
\tilde a_{it} &\sim Q\big(\,\cdot \mid \tilde y_{it}, o_{it}, \tilde\pi_{it}\big),
\qquad \tilde \pi_{i,t+1} = U(\tilde\pi_{it}, \tilde y_{it}, \tilde a_{it}),\\
\tilde x_{i,t+1} &= \tilde x_{it} \text{ with } (\tilde y_{it}, \tilde a_{it})
\text{ appended as the previous decision and the new acquisitions.}
\end{align}
After the first holdout occasion nothing in the decision or acquisition streams
is real. The model conditions on what \emph{it} decided and on what the lottery
gave it. Inventory, duplicate counts and every satiation term update
automatically, because they are functions of the acquisition stream.

\subsection{Teacher forcing}

A decision is drawn at every occasion, but it is recorded and never fed back:
\begin{equation}
\tilde y_{it} \sim p_\theta\big(\,\cdot \mid x_{it}^{\text{obs}}\big),
\qquad t \in \mathcal{H}_i .
\end{equation}
Each draw is a one-step forecast from the true history. The scored history is
$(\tilde y, a^{\text{obs}})$.

\subsection{Arm C: real decisions, lottery re-drawn}

No model is involved. The customer's observed decisions are replayed, and the
environment re-draws what they yielded:
$\tilde a_{it} \sim Q(\,\cdot \mid y_{it}^{\text{obs}}, o_{it}, \tilde\pi_{it})$.
This isolates the randomness of the loot boxes, together with any error in the
environment.

\subsection{The two-by-two}

The four arms differ only in which streams are fed back:

\begin{center}
\small
\begin{tabular}{lll}
\toprule
 & acquisitions real & acquisitions simulated \\
\midrule
decisions real    & the data itself & arm C \\
decisions sampled & teacher forcing & free running \\
\bottomrule
\end{tabular}
\end{center}

\paragraph{Only free running is a simulator.} Arm C replays each customer's
actual holdout decisions. Teacher forcing feeds the model the customer's actual
holdout decisions and acquisitions at every occasion. Both therefore need the
realised future, which by definition does not exist for a policy that has not
been run. They have no managerial use as simulators. Their role is diagnostic:
each switches off one feedback channel, which locates where the failure enters.
The yardstick for a generator a manager could actually use is free running,
judged against the baselines below. The results table keeps the two groups
apart, and they should not be ranked on one scale.

\subsection{Baselines}

Both baselines use the same grid, environment, warm-up and replicates as the
model, so a difference between them and the model is a difference in the
decision process alone. Both are estimated on the same training customers'
calibration rows as the environment.
\begin{itemize}
\item \textbf{Independent draws:} $\tilde y_{it} \sim \hat q(y)$, the
  calibration-period shares of the nine decisions. This matches every marginal
  share by construction and must fail anything that depends on order.
\item \textbf{First-order Markov chain:}
  $\tilde y_{it} \sim \hat q(y \mid \tilde y_{i,t-1})$, a $9\times 9$
  transition matrix plus a row for ``no previous decision'', Laplace-smoothed
  ($\alpha = 1$). This 81-number model is the bar a deep sequence model has to
  clear.
\end{itemize}

\subsection{Scheduled sampling}

The rollout is identical to free running. What changes is how the model was
\emph{trained} \citep{bengio2015}. In each minibatch:
\begin{enumerate}
\item A first forward pass under teacher forcing draws
  $\hat y_{t} \sim p_\theta(\cdot \mid x_t^{\text{obs}})$ at every row.
\item Each row's previous-decision input is replaced by the model's own draw
  $\hat y_{t-1}$ with probability $p_e = p\cdot\min\{1, (e+1)/10\}$, which ramps
  to $p$ over the first ten epochs $e$. Otherwise it keeps the observed
  $y_{t-1}$.
\item The loss is computed on a second pass with the corrupted inputs.
\end{enumerate}
Only the decision stream is corrupted; acquisitions stay real in training. The
objective is known to be inconsistent \citep{huszar2015}, so it is reported as
a remedy, not as a contribution.

\section{Measuring the distance between sets of histories}\label{sec:metric}

\subsection{Step 1: summarise each history as 27 statistics}

Each holdout history $h$, real or generated, is reduced to a vector
$\phi(h) \in \mathbb{R}^{27}$. The model was not fitted to any of these
statistics. All are computed over the holdout occasions only.

\begin{center}
\small
\begin{tabularx}{\linewidth}{@{}l c X@{}}
\toprule
statistic & count & definition \\
\midrule
spend & 1 & total pulls bought (Buy1 $=1$, Buy10 $=10$) \\
incidence & 1 & share of occasions with a purchase \\
decision shares & 9 & share of occasions with each of the nine decisions \\
longest NotBuy run & 1 & longest run of consecutive NotBuy occasions \\
buy runs & 2 & mean and longest run of consecutive purchase occasions \\
switch rate & 1 & share of consecutive purchases that change banner \\
spend concentration & 1 & Herfindahl index of spend across the four banners \\
autocorrelation & 5 & of the purchase indicator, lags 1--5 \\
escalation & 1 & $P(\text{Buy10 at } t \mid \text{Buy10 at } t-1)$ \\
persistence & 1 & $P(\text{purchase at } t \mid \text{purchase at } t-1)$ \\
satiation signature & 4 & $P(\text{buy on banner } b \mid \text{holding } k
  \text{ copies of } b\text{'s featured product})$, $k = 0,1,2,3{+}$, pooled over
  figure A, figure B and weapon \\
\bottomrule
\end{tabularx}
\end{center}

\noindent The satiation signature is the statistic the inventory stream exists
to get right. It is computed from the acquisition stream that accompanies the
decisions: the real one for real and teacher-forced histories, the simulated
one for free running, arm C and the baselines.

\subsection{Step 2: standardise}

Each statistic is centred and scaled by the \emph{real} set's mean and standard
deviation. Undefined values (escalation for a customer who never bought a
Buy10, an autocorrelation with too few occasions, an empty satiation bucket)
are set to the real mean, that is to zero after standardisation, for both sets
alike. A missing value cannot by itself create a discrepancy. Distances below
are therefore measured in units of real-data standard deviations, in a
27-dimensional space.

\subsection{Step 3: energy distance}

Let $X$ and $X'$ be independent draws of $\phi(h)$ from the real distribution
$F$, and $Y, Y'$ independent draws from the generated distribution $G$. The
energy distance \citep{szekely2013} is
\begin{equation}
\mathcal{E}(F,G) \;=\; 2\,\E\norm{X - Y} \;-\; \E\norm{X - X'} \;-\; \E\norm{Y - Y'},
\label{eq:ed}
\end{equation}
with $\norm{\cdot}$ the Euclidean norm. In words: twice the average distance
between a real and a generated history, minus the average distance among real
histories and among generated ones. If the generator is faithful, a generated
history sits as far from a real one as two real ones sit from each other, and
the three terms cancel.

Two properties make it the right yardstick here. First,
$\mathcal{E}(F,G) \ge 0$, with equality if and only if $F = G$ (given finite
first moments). Second, it compares the whole joint distribution, not just
means, because it equals a weighted distance between characteristic functions,
\begin{equation}
\mathcal{E}(F,G) \;=\; \frac{1}{c_d}\int_{\mathbb{R}^d}
\frac{|\varphi_F(s) - \varphi_G(s)|^2}{\norm{s}^{d+1}}\, ds,
\qquad c_d = \frac{\pi^{(d+1)/2}}{\Gamma\big((d+1)/2\big)},
\end{equation}
with $d = 27$. A difference in level, spread, shape or co-movement of the
statistics all count.

\paragraph{The estimator is signed.} With samples $x_1,\dots,x_n$ and
$y_1,\dots,y_m$,
\begin{equation}
\widehat{\mathcal{E}} \;=\; \frac{2}{nm}\sum_{i=1}^{n}\sum_{j=1}^{m}\norm{x_i - y_j}
\;-\; \frac{1}{n(n-1)}\sum_{i \ne i'}\norm{x_i - x_{i'}}
\;-\; \frac{1}{m(m-1)}\sum_{j \ne j'}\norm{y_j - y_{j'}} .
\label{eq:edhat}
\end{equation}
The within-set sums leave out each point's zero distance to itself, which makes
$\widehat{\mathcal{E}}$ unbiased: $\E\,\widehat{\mathcal{E}} = \mathcal{E}(F,G)$
at every sample size. When $F = G$ it therefore scatters around zero and is
negative about half the time. A negative value, such as arm C's $-0.010$, means
``indistinguishable from real'', not ``closer than real''.

\subsection{Step 4: the reference band}

A raw distance has no natural scale, so it is calibrated against real data
scored against itself. Let $R$ be the $n_R$ real customers, $S$ the
$n_R \cdot M$ generated histories, and $h = \lfloor n_R/2 \rfloor$.

\begin{enumerate}
\item \textbf{Null distribution.} For $b = 1,\dots,B$, split $R$ at random into
  two disjoint halves of size $h$ and compute $\widehat{\mathcal{E}}_b^{0}$
  between them. With $n_R = 300$, $h = 150$ and $B = 120$ (R36i, as in R36c
  and R36e).
\item \textbf{Observed value.} For $b = 1,\dots,B$, draw $h$ customers from $R$
  and $h$ histories from $S$, each without replacement, compute
  $\widehat{\mathcal{E}}_b$, and average:
  $\bar{\mathcal{E}} = B^{-1}\sum_b \widehat{\mathcal{E}}_b$. Drawing at the
  same size $h$ makes each draw directly comparable with a null draw.
\item \textbf{Verdict.} The generator is \emph{inside the band} if
  $\bar{\mathcal{E}} \le q_{0.95}\big(\{\widehat{\mathcal{E}}_b^{0}\}\big)$.
  The scripts also report
  $z = \big(\bar{\mathcal{E}} - \text{mean}(\widehat{\mathcal{E}}^0)\big) /
  \text{sd}(\widehat{\mathcal{E}}^0)$, and the share of null draws at or above
  $\bar{\mathcal{E}}$.
\end{enumerate}

\subsection{How to read the verdict}

The averaged statistic $\bar{\mathcal{E}}$ estimates the population distance
$\mathcal{E}(F,G)$, and its expectation does not depend on $h$. The band does:
under $F=G$ the estimator's spread shrinks roughly as $1/h$. The recorded bands
fit that rate: 0.373 at 30 against 30, and 0.078 at 150 against 150, a ratio of
4.8 against a sample-size ratio of 5.

Two consequences follow.
\begin{itemize}
\item \textbf{It is a tolerance criterion, not a 5\% test.} The rule compares an
  average over many draws with the spread of a single draw. It asks whether the
  generator's distance from the real data is no larger than the distance between
  two random halves of real customers typically is. A generator with a small but
  real discrepancy ($0 < \mathcal{E} < 0.078$ at this sample size) passes. That
  is a defensible standard, but the paper should state it in these words, or
  replace it with a permutation test on the full samples if a formal test is
  wanted.
\item \textbf{Verdicts at different sample sizes are not comparable.} A distance
  that passes at 60 customers can fail at 300, because the band narrows while
  the estimate does not. This is exactly what happened to the teacher-forced
  arm (Section~\ref{sec:caveats}).
\end{itemize}

\subsection{The weakness: dependence among the statistics}\label{sec:weakness}

Energy distance reacts strongly when the distribution of a statistic is wrong,
and weakly when every statistic's distribution is right but the statistics are
wrongly \emph{linked} across customers \citep{pinson2013}. In marketing terms, a
generator could reproduce the spread of spend across customers and the spread of
satiation across customers, yet attach strong satiation to light spenders
instead of heavy ones. It would be generating the wrong customer types, which
is what targeting depends on. Energy distance would barely register it.

\paragraph{Why.} Take the squared distance between a real and a generated
history. With $X$ and $Y$ independent, and means $\mu$ and $\nu$,
\begin{equation}
\E\norm{X - Y}^2 \;=\; \sum_{k=1}^{d}\Big[\operatorname{Var}X_k +
\operatorname{Var}Y_k + (\mu_k - \nu_k)^2\Big].
\end{equation}
No covariance appears: the squared distance depends on each statistic's
distribution separately. Energy distance uses the unsquared distance, so the
dependence structure enters only through how spread out $\norm{X-Y}^2$ is
around its mean, which is a second-order effect. In 27 dimensions the distance
also concentrates tightly around its root mean square, which shrinks that
second-order effect further. A delta-method expansion makes this exact enough to
quantify. For Gaussian statistics with unit variances, dimension $d$ and
covariance matrices $\Sigma_X, \Sigma_Y$:
\begin{align}
\mathcal{E} &\;\approx\; \frac{\norm{\Sigma_X - \Sigma_Y}_F^2}{2\,(2d)^{3/2}}
&&\text{margins right, dependence wrong,} \label{eq:weak}\\
\mathcal{E} &\;\approx\; \frac{\norm{\mu - \nu}^2}{\sqrt{2d}}
&&\text{dependence right, means wrong.} \nonumber
\end{align}
With $d = 27$ the two denominators are 794 and 7.3. Per unit of squared error, a
wrong mean weighs about 110 times as much as a wrong covariance.

\paragraph{How large, in practice.} Table~\ref{tab:weak} simulates 27 standard
normal statistics at the R36 sample sizes (300 real customers against 2,400
generated histories, 120 split-half draws), using the same code as R36.

\begin{table}[h]
\centering
\small
\caption{Energy distance against a wrong dependence structure and against a
wrong mean, 27 standard normal statistics. ``Predicted'' is
equation~\eqref{eq:weak}; ``population'' uses 2,500 against 2,500, three
repetitions; ``R36 procedure'' is \texttt{seq\_metrics.compare} at R36's sizes,
with that case's own band. The classifier is trained on 300 against 300.}
\label{tab:weak}
\begin{tabular}{@{}lrrrrcr@{}}
\toprule
what the generator gets wrong & predicted & population & R36 proc. & band & verdict & AUC \\
\midrule
nothing & 0.000 & 0.001 & $-0.001$ & 0.024 & inside & 0.45 \\
one correlation, $+0.5 \to -0.5$ & 0.003 & 0.002 & 0.004 & 0.023 & inside & 0.59 \\
all 351 correlations, $0.3 \to 0$ & 0.080 & 0.050 & 0.048 & 0.048 & at the edge & 0.78 \\
all 351 correlations, $0.5 \to 0$ & 0.221 & 0.127 & 0.113 & 0.046 & outside & 0.91 \\
one mean, shifted by 0.75 s.d. & 0.076 & 0.078 & 0.079 & 0.021 & outside & 0.62 \\
\bottomrule
\end{tabular}
\end{table}

\noindent Reversing the sign of a correlation, a qualitative error about which
customers do what, moves the distance by 0.004, a sixth of its band. Deleting
every one of the 351 correlations of 0.3 only just reaches the band. Shifting a
single statistic's mean by three quarters of a standard deviation lands at
almost four times its band. The expansion is accurate for small errors and
overstates large ones, but the ordering is not in doubt. The bands here are for
Gaussian statistics; the real band (0.078) differs because the real statistics
are skewed, discrete and partly imputed, so each case should be read against
its own band.

\paragraph{The consequence for the figure.} A generator inside the band is
shown to have the right distribution for each statistic and roughly the right
overall spread. It is not shown to have the right links between statistics.
Note also that the classifier does the opposite: it detects the correlation
errors better than the mean shift. The per-statistic tables in
Appendix~\ref{app:perstat} cannot help here at all: a one-dimensional distance
is blind to dependence by construction. A direct check would compare the
$27\times27$ correlation matrices of real and generated statistics, with a
split-half band for that comparison too.

\paragraph{The complements.} Because dependence is what a rollout puts at risk,
the pipeline reports two further measures.
\begin{itemize}
\item A \textbf{classifier two-sample test} \citep{lopezpaz2017}: gradient-boosted
  trees trained to tell real from generated, five-fold cross-validated AUC with
  a 95\% interval. An AUC of 0.5 means indistinguishable. Permutation importances
  name the statistics the classifier leans on, which is the diagnosis.
\item A \textbf{variogram score} \citep{scheuerer2015}, which compares the
  spread of every pair of statistics and is sensitive to dependence.
\end{itemize}

\section{The numbers on the fixed engine (R36i)}\label{sec:results}

R36i (job 42879, 20 October 2026) scored every arm in one process: one draw of
300 customers, eight replicates, one scaler and one band, on the rollout engine
fixed on 16 October. The band reproduced the earlier runs' $q_{0.95}$ of
0.0781258 exactly, so these are the same customers as before. Every arm run
generated a decision on exactly the 179,920 live occasions ($8 \times 22{,}490$),
with no mismatched occasions and no real acquisition altered.

\begin{table}[h]
\centering
\small
\caption{Energy distance between generated and real holdout histories, R36i.
Mean $\pm$ s.d.\ over three runs; band $q_{0.95} = 0.078$; AUC is the mean
classifier two-sample AUC. The last column is the pre-fix value.
$^{*}$Pre-fix teacher forcing was b29\_hl2 on 60 customers, against a band of
0.373.}
\label{tab:results}
\begin{tabular}{@{}lrrrrr@{}}
\toprule
arm & $\bar{\mathcal{E}}$ & $\times$ band & $z$ & AUC & pre-fix \\
\midrule
\multicolumn{6}{@{}l}{\emph{Generators a manager could use}}\\
first-order Markov chain & $0.784 \pm 0.006$ & 10.0 & 21.8 & 0.885 & 0.711 \\
scheduled sampling, $p = 0.35$ (b35) & $1.311 \pm 0.299$ & 16.8 & 36.5 & 0.841 & 0.928 \\
scheduled sampling, $p = 0.15$ (b35) & $1.710 \pm 0.389$ & 21.9 & 47.5 & 0.851 & 1.390 \\
independent draws & $2.468 \pm 0.008$ & 31.6 & 68.6 & 0.913 & 2.190 \\
free running (b33) & $3.169 \pm 0.381$ & 40.6 & 88.1 & 0.869 & 2.352 \\
\midrule
\multicolumn{6}{@{}l}{\emph{Diagnostic controls: they read the realised holdout}}\\
arm C & $-0.000 \pm 0.002$ & 0.0 & 0.0 & 0.552 & $-0.010$ \\
teacher forcing (b33) & $0.417 \pm 0.080$ & 5.3 & 11.6 & 0.851 & $0.307^{*}$ \\
\bottomrule
\end{tabular}
\end{table}

\subsection{The pre-registered verdicts}

\begin{center}
\small
\begin{tabular}{@{}lp{0.56\linewidth}@{}}
\toprule
verdict & result \\
\midrule
V1 the lottery is harmless & holds: arm C inside the band on 3 of 3 draws \\
V2 the closed loop fails & holds: free running outside on 3 of 3 seeds \\
V3 the per-step law is indistinguishable & \textbf{fails}: teacher forcing outside on 3 of 3 seeds \\
V4 the loop makes most of the gap & holds: teacher forcing removes 88\%, 85\% and 87\% of the free-running distance \\
V5 the ordering & holds: arm C $<$ teacher forcing $<$ Markov $<$ free running \\
\bottomrule
\end{tabular}
\end{center}

\noindent The diagnosis therefore survives in its weak form, as pre-registered.
The paper can say that \emph{most of the distance is created by feeding the
model its own output}. It can no longer say that teacher-forced sequences are
indistinguishable from real ones.

\paragraph{What goes wrong.} Pooled over the holdout occasions, free running
abandons NotBuy and inflates the rare multi-pull decisions:

\begin{center}
\small
\begin{tabular}{@{}lrrrrr@{}}
\toprule
 & NotBuy & Buy10 regular & Buy10 figure B & Buy1 weapon & Buy10 weapon \\
\midrule
observed & 0.425 & 0.003 & 0.013 & 0.044 & 0.010 \\
teacher forcing & 0.352 & 0.012 & 0.032 & 0.061 & 0.020 \\
free running & 0.156 & 0.033 & 0.111 & 0.159 & 0.183 \\
\bottomrule
\end{tabular}
\end{center}

\paragraph{The remedy.} Energy falls monotonically in the scheduled-sampling
probability: 3.169, 1.710, 1.311 for $p = 0, 0.15, 0.35$. Both remedied models
are lower than b33 on every seed. At $p = 0.35$ the remedy closes 59\% of the
gap (61\% before the fix), but it is still 17 times the band and behind the
Markov chain on every seed (its best seed 1.129, against 0.791).

\subsection{Which statistics give each generator away}\label{sec:perstat}

Table~\ref{tab:overview} applies the same yardstick to each statistic alone
(Appendix~\ref{app:perstat} has the definition and the 27 full tables). Each
entry is the arm's mean energy as a multiple of that statistic's own band.

\begin{table}[h]
\centering
\small
\caption{Energy by statistic, in multiples of each statistic's real-vs-real
band (R36i, mean over three runs). Grey: inside the band on every run. Bottom
row: statistics outside the band on every run.}
\label{tab:overview}
\InputIfFileExists{r36i_per_statistic_overview.tex}{}{\emph{(overview not
generated)}}
\end{table}

\begin{enumerate}
\item \textbf{Free running fails almost everywhere, mainly by overspending.} It
  is outside on every run for 24 of 27 statistics. Customers buy 349 pulls on
  average over the holdout against 110, choose NotBuy on 9\% of occasions
  against 28\% (averaged per customer), and their longest NotBuy run is 3.9
  occasions against 10.9. The worst statistic is the Buy10 weapon share, 179
  times its band.
\item \textbf{Scheduled sampling halves the damage and keeps the pattern.} At
  $p = 0.35$, 15 of 27 statistics are outside on every run, and spend is still
  double (219 pulls). It does repair the banner switch rate and spend
  concentration.
\item \textbf{The Markov chain gets the levels right and misses longer memory.}
  Spend, the mean buying run, persistence and escalation are inside on every
  run. It fails on everything that needs more than one step of history: the
  purchase autocorrelation at lags 2--5 (30 to 42 times the band) and the
  longest NotBuy run (16 times).
\item \textbf{Independent draws fail everything that depends on order}: the
  switch rate is 79 times its band and the lag-1 autocorrelation 101 times.
\item \textbf{Arm C passes 26 of 27.} Its one failure is the satiation signature
  at one copy (7.7 times the band). The lottery leaves 68\% of customers without
  a defined value, against 50\% in the data, and those who have one buy again at
  0.26 against 0.12. The Markov chain (7.4) and independent draws (6.9) fail
  this statistic by the same amount, so for them it is the lottery's error. For
  free running (58 times) most of it is the model's.
\item \textbf{Teacher forcing misses run lengths and switching}: 16 of 27 are
  outside on every run, the worst being the lag-1 autocorrelation (17 times),
  the switch rate (16) and the longest NotBuy run (9.2). This matches the joint
  classifier, which separates teacher forcing from real data almost as well as
  it separates free running (AUC 0.851 against 0.869).
\end{enumerate}

\subsection{The same yardstick in later work}

Two later experiments reused this yardstick, with the same customers, band and
lottery, and reproduced R36i's model-free arms to the fourth decimal or better.
\begin{itemize}
\item \textbf{R43 step 2 (29 October).} The gen-6 streaming model and the tuned
  gen-5 model free-run at 2.617 and 2.563 (34 and 33 times the band), with no
  difference between them. Teacher forcing removes 76--90\% of the distance in
  both, and the Markov chain is ahead of every seed of every family.
\item \textbf{R46 part B (27 October; branch \texttt{r46-parametric}, not yet
  merged).} Classical parametric choice models with lagged decisions, run
  structure and holdings free-run at 0.308 (R4) and 0.318 (R6, with latent
  classes): ten times closer than free-running ProductGPT and closer than the
  Markov chain, although R6 predicts worse one step ahead (holdout NLL 0.983
  against the GRU's 0.886). Nothing passes the band. This is the sharpest instance of the
  paper's spine: the weaker predictor is by far the better simulator.
\end{itemize}

\section{Qualifications}\label{sec:caveats}

\paragraph{1. The earlier numbers ran on a rollout engine with four defects.}
They were found and fixed on 16 October 2026 (commit \texttt{7aca2a3}), and
R36i re-ran every arm on the fixed engine:
\begin{enumerate}
\item \textbf{Pre-window inventory dropped.} The rollout never passed the
  inventory accumulated before the 384-row window, so slots-type models
  (b29\_hl2, b33, b35) started every rollout with an empty inventory.
\item \textbf{Wrong item mix.} The environment drew every 4-star from the
  character pool, even on the weapon banner, and never gave a regular-banner
  5-star weapon. The real 4-star weapon shares are 0.32--0.74 by banner; 5-star
  weapons are 0.23 of regular-banner 5-stars.
\item \textbf{Early loop exit.} The loop stopped at the first row where no
  customer in a batch was active, so a customer whose holdout began after such
  a gap was scored as an empty history.
\item \textbf{Pity warm-up off by one.} The warm-up stopped at the batch's
  earliest holdout row and missed each customer's last real outcome.
\end{enumerate}

\noindent \emph{What the fix changed.} Free running rose from 2.352 to 3.169
(by 0.62 to 1.00 per seed) and now loses to independent draws on every seed.
The arms without a model moved too (independent draws by $+0.28$, Markov by
$+0.07$). Only defects 3 and 4 can reach them, so part of every pre-fix
number's error was the engine's, shared by all arms. Free running rose about
three times as much as independent draws, which is consistent with defects 1
and 2 (what the model feeds itself) adding to it; this run cannot separate the
four. Defect 3 did reach the pre-fix runs: it skipped 939 of the 22,490 holdout
occasions (4.2\%), belonging to 10 of the 300 customers. Still on the pre-fix
engine, and not to be quoted without that qualifier: the R36c table comparing
model variants (b18 cells, b29\_hl2), R37 (the duplet), R38 and R42.

\paragraph{2. Teacher forcing is not indistinguishable from real data.} Before
the fix its verdict ``inside the band'' compared 0.307 with a 60-customer band
of 0.373, five times wider than the 300-customer band
(Section~\ref{sec:metric}). On 300 customers it is 0.417, 5.3 times the band,
outside on every seed, at $z = 11.6$. It is also \emph{detectable}: the
classifier separates it from real data almost as well as free running (AUC
0.851 against 0.869), on run lengths and switching. It generates NotBuy on
35\% of occasions against 43\% observed. What survives is the contrast: it
removes 85--88\% of the free-running distance. The defensible statement is
that \emph{most} of the generative gap comes from feeding the model its own
output, not that the per-step distribution is exact.

\paragraph{3. ``Inside the band'' is a tolerance, and it ignores dependence.}
The rule asks whether a generator lands no farther from the data than two random
halves of real customers do; it is not a 5\% test (Section~\ref{sec:metric}).
Being inside also says nothing about whether the statistics are correctly linked
across customers, because energy distance barely sees dependence
(Section~\ref{sec:weakness}).

\paragraph{4. The lottery has a small, located error.} Arm C's classifier
interval excludes 0.5 on every draw (0.518--0.607), so under R36's stricter rule
(the classifier's interval must contain 0.5) not even arm C passes. The
per-statistic breakdown locates the error: the satiation signature at one copy,
7.7 times its band, and nothing else. Every arm that draws acquisitions from the
lottery inherits it. The identity process is also calibrated rather than taken
from the game, because the data's featured shares are incompatible with the
game's guarantee. The lottery has to be tightened before it carries a policy
number.

\paragraph{5. What the design cannot test.} The occasion grid is fixed, so the
exercise says nothing about purchase timing or frequency; that needs a duration
head. The 384-occasion window is priced at $+0.0007$ nats for b33, with the
pre-window inventory carried in.

\paragraph{Status.} Every number in Section~\ref{sec:results} is from R36i
(commit \texttt{3ef735e} on \texttt{gen6-streaming}) and can be quoted without
the pre-fix qualifier, for the gen-5 b33 and b35 models at context 384. The
figure now uses these numbers; arm C and teacher forcing appear only in its
two-by-two, as diagnostic controls, and its yardstick panel ranks only the
generators a manager could use.

\appendix
\section{Energy distance one statistic at a time}\label{app:perstat}

The joint distance says how far a generator lands; the tables below say on which
statistics. For statistic $k$, with $F_k$ and $G_k$ its real and generated
distribution functions,
\begin{equation}
\mathcal{E}_k \;=\; 2\,\E|X_k - Y_k| - \E|X_k - X_k'| - \E|Y_k - Y_k'|
\;=\; 2\int_{-\infty}^{\infty}\big(F_k(x) - G_k(x)\big)^2\,dx,
\end{equation}
twice the squared area between the two distribution functions. Each table uses
the same standardised values, the same rule for undefined values, the same
estimator and the same split-half draws as the joint distance, so each
statistic has its own band built on the same customer splits. The source is
R36i's saved per-customer statistics, which reproduce all 21 run-level joint
energies behind Table~\ref{tab:results} to within $10^{-15}$. The tables are generated by
\texttt{scripts/r36i\_per\_statistic.py}.

\paragraph{Columns.} \emph{Energy} is the mean $\pm$ s.d.\ over the three runs
of each arm (training seeds for model arms, draw seeds otherwise). \emph{Inside}
counts the runs at or below the band. \emph{Gen.\ mean} is the generated
statistic's average in its natural units, to set against the real mean in the
caption. \emph{Undefined} is the share of generated histories for which the
statistic does not exist (no Buy10 to escalate from, an empty satiation
bucket).

\paragraph{Four cautions.}
\begin{enumerate}
\item The tables are blind to dependence by construction
  (Section~\ref{sec:weakness}). A generator can pass all 27 and still fail
  jointly.
\item Undefined values are set to the real mean, as in the joint distance. A
  generator that leaves a statistic undefined more often than real customers do
  therefore shows a distance on it. That is informative in itself, which is why
  the undefined share is shown. For the satiation signature at $k = 2$ and
  $k = 3{+}$, undefined for 89\% and 97\% of real customers, the bands are
  tiny (0.003 and 0.001), and the ratios mostly reflect how often the statistic
  is defined.
\item With 27 statistics and a 95\% band, an exact generator would still land
  outside on about one statistic in twenty. Read the pattern across statistics,
  not a single exceedance.
\item The lottery's own error sits on the satiation signature at one copy
  (arm C, 7.7 times its band), and every arm that uses the lottery inherits it.
\end{enumerate}

\InputIfFileExists{r36i_per_statistic_tables.tex}{}{%
\begin{center}\fbox{\parbox{0.9\linewidth}{\small No tables yet. They are
generated by \texttt{scripts/r36i\_per\_statistic.py} from the per-customer
statistics that the R36i run saves. The pre-fix runs saved only summaries, so
their per-statistic distances cannot be recovered without re-running them.}}
\end{center}}

\begin{thebibliography}{9}
\bibitem[Bengio et al.(2015)]{bengio2015}
Bengio, S., Vinyals, O., Jaitly, N., and Shazeer, N. (2015).
Scheduled sampling for sequence prediction with recurrent neural networks.
\emph{Advances in Neural Information Processing Systems} 28.

\bibitem[Husz\'ar(2015)]{huszar2015}
Husz\'ar, F. (2015). How (not) to train your generative model: scheduled
sampling, likelihood, adversary? arXiv:1511.05101.

\bibitem[Lopez-Paz and Oquab(2017)]{lopezpaz2017}
Lopez-Paz, D. and Oquab, M. (2017). Revisiting classifier two-sample tests.
\emph{International Conference on Learning Representations}.

\bibitem[Pinson and Tastu(2013)]{pinson2013}
Pinson, P. and Tastu, J. (2013). Discrimination ability of the energy score.
Technical report, Technical University of Denmark.

\bibitem[Scheuerer and Hamill(2015)]{scheuerer2015}
Scheuerer, M. and Hamill, T. M. (2015). Variogram-based proper scoring rules for
probabilistic forecasts of multivariate quantities. \emph{Monthly Weather
Review} 143, 1321--1334.

\bibitem[Sz\'ekely and Rizzo(2013)]{szekely2013}
Sz\'ekely, G. J. and Rizzo, M. L. (2013). Energy statistics: A class of
statistics based on distances. \emph{Journal of Statistical Planning and
Inference} 143(8), 1249--1272.
\end{thebibliography}

\end{document}
